% Файл намеренно идемпотентен: повторное подключение ничего не печатает.
\ifcsname hiddenpowerabstractprinted\endcsname
\else
  \expandafter\gdef\csname hiddenpowerabstractprinted\endcsname{1}

  \chapter*{Аннотация}
  \addcontentsline{toc}{chapter}{Аннотация}

  \begingroup
  \setlength{\emergencystretch}{3em}%
  \begin{sloppypar}
  Настоящий конспект систематизирует математический аппарат, необходимый для исследовательской работы в области машинного обучения и генеративного проектирования. Основное внимание уделяется не изолированному изучению отдельных дисциплин, а связям между линейной алгеброй, многомерным анализом, дифференциальной геометрией, динамическими системами, обыкновенными и стохастическими дифференциальными уравнениями, уравнениями в частных производных и численными методами.

  Линейная алгебра излагается с опорой на курс Э.\,Б.~Винберга и дополняется материалом, необходимым для современных вычислительных задач: спектральными разложениями, сингулярным разложением, псевдообратными матрицами, матричным дифференцированием и тензорными конструкциями. Последующие разделы показывают, как эти понятия переходят в геометрию многомерных пространств, анализ непрерывной динамики, диффузионные и потоковые генеративные модели, физические симуляции, обратные задачи и оптимизацию формы.

  Изложение строится по единой схеме: точное определение, содержательное объяснение, основные утверждения, необходимые элементы доказательства, вычислительный пример и связь с другими разделами. Строгость сохраняется в формулировках и условиях применимости результатов, однако технические доказательства приводятся только тогда, когда они необходимы для понимания метода.

  Конспект не претендует на полноту университетских курсов по перечисленным дисциплинам. Его задача состоит в построении целостной математической основы, позволяющей читать современные исследовательские работы, понимать используемые модели и самостоятельно реализовывать методы генеративного проектирования.
  \end{sloppypar}
  \endgroup
\fi
